\section{Discussion}
\label{sec:discussion}

\para{A single mechanism explains all three experiments.} The probe (\secref{sec:res2}) is
the hinge. The base model encodes its answer in the prompt representation before it generates
($0.87$ decodable). This one fact predicts its behavior everywhere else. It is rigid: because
the answer is already fixed by the prompt, feedback \emph{appended after} the answer---whether
empty pushback ($0.17$ change) or a genuine hint ($0.02$ change)---cannot move it
(\secref{sec:res1}). Yet its uncertainty is highly shaped by \emph{what is in the prompt}:
placing a fact \emph{before} the answer redistributes its entropy by $0.32$ bits
(\secref{sec:res3}). Both behaviors follow from the same cause---the answer is a function of
the prompt representation. Information in the prompt shapes the commitment; information after
it does not. This resolves what looks like a contradiction between Experiment~1 (base ignores
a hint) and Experiment~3 (base is strongly moved by a fact): the difference is whether the
information arrives before or after the model has committed.

\para{Post-training moves commitment later.} Instruct decides during generation, which makes
it reactive---but plain RLHF makes it \emph{indiscriminately} so. It over-reacts, abandoning
correct answers under content-free challenge ($\Delta^{c\rightarrow i}$ up to $0.38$). This is
the sycophancy/FlipFlop failure~\citep{laban2023flipflop}, the opposite of pre-commitment. The
reasoning model commits latest of all (probe $0.33$, at chance) and is both the most selective
reactor and the most robust in its uncertainty: it reasons over the whole question rather than
reading its answer off the prompt. On the hypothesis's own terms, it ``knows the answer it
wants to give'' the \emph{least}.

\para{Reconciling the contested direction.} Our \emph{internal} mechanism is clean and
monotonic, but the \emph{behavioral} reactivity direction among post-trained models depends on
the task and on whether the model was right to begin with. On \arc, where models are often
wrong and the hint is informative, the reasoning model reacts beneficially; on \gsm, where it
is usually right and feedback is content-free, the same model over-reacts and breaks correct
solutions. What is invariant across tasks is \emph{base-model rigidity}. This is exactly the
``contested direction'' the literature flagged~\citep{ufo2025,kumar2024score,laban2023flipflop}:
under- and over-reaction coexist and are selected by conditions. We therefore report both
rather than claiming a single direction, and we locate the two classic failure modes at
opposite ends of one post-training axis---UFO/SCoRe-style under-reaction in the \emph{base}
model, FlipFlop-style over-reaction in \emph{instruct}.

\para{Limitations.} We are honest about the scope of these claims.
\begin{itemize}[leftmargin=*,itemsep=1pt,topsep=2pt]
    \item \textbf{One family, one scale (1.5B).} Lineage is matched, but \instruct and \rlm
    differ in the \emph{type} of post-training (RLHF vs reasoning distillation), not a clean RL
    ablation. The result should be confirmed at 7B and on other families (\eg SimpleRL-Zoo,
    Qwen-UFO).
    \item \textbf{Interface differs by model} (base$=$transcript, others$=$chat), because base
    models have no chat mode. This is inherent to the comparison but a confound for the
    \emph{absolute} reactivity rates in Experiment~1. The internal probe (Experiment~2),
    measured identically across models, is the cleaner contrast---and it agrees.
    \item \textbf{Low \arc accuracy for \rlm (37\%)} reflects a math-distilled model on science
    multiple-choice under a 640-token budget (reasoning is sometimes truncated). Because the
    probe targets the model's own answer, this does not bias Experiment~2, but turn-1 accuracy
    should be read with care.
    \item \textbf{Coarse entropy.} Experiment~3 estimates entropy from $K=6$ samples per item;
    we rely on 120-item aggregates and paired tests, which are robust, and replaced a
    noise-inflated per-item $|\Delta H|$ metric with signed bootstrap-CI reductions.
    \item \textbf{Heuristics.} Greedy decoding with a fixed token budget and sentence-level
    \strat fact-splitting are heuristics; parse-failure rates ($\leq6\%$, mostly the reasoning
    model) are logged and reported, never silently dropped.
\end{itemize}

\para{Broader implications.} The ``pre-plans and cannot react'' property that motivates
safety concern is real, but it is a property of \emph{base/pretrained} models and of
information that arrives \emph{after} the model has committed. RL post-training largely
\emph{mitigates} it---at the cost, for plain RLHF, of the opposite failure. Practically, this
suggests the right lever for mid-conversation correctability is not avoiding RL but choosing a
post-training recipe that decides late \emph{and} discriminately, as the reasoning model does
on \arc. It also suggests rotating certainty could become a decoding-time probe for residual
per-hop uncertainty---useful precisely where a model's answer is prompt-driven.
